% =============================================================================
%  Project Proposal -- os-buddy
%  Cover and design follow the "Cubicle - Project Report" (Times text, plain
%  black headings, grey-header tables, DIU cover page).
%  Build:  pdflatex proposal.tex   (run twice so cross-references resolve)
%  Edit the "Project details" and "Team" blocks below; everything else follows.
% =============================================================================
\documentclass[11pt,a4paper]{article}

% ---------- Fonts: Times text, Latin Modern Mono for code --------------------
\usepackage[T1]{fontenc}
\usepackage[utf8]{inputenc}
\usepackage{lmodern}                 % provides the Latin Modern Mono code font
\usepackage{times}                   % Times (Nimbus Roman) for text
\renewcommand{\ttdefault}{lmtt}
\usepackage{microtype}

% ---------- Layout ----------------------------------------------------------
\usepackage[a4paper,top=2.3cm,bottom=2.2cm,left=2.5cm,right=2.5cm,headheight=14pt]{geometry}
\usepackage{fancyhdr}
\usepackage{graphicx}
\usepackage{xcolor}
\usepackage{array}
\usepackage{tabularx}
\usepackage{caption}
\usepackage[hidelinks]{hyperref}

% ---------- Project details (edit these) ------------------------------------
\newcommand{\projectname}{os-buddy}
\newcommand{\projectsubtitle}{An Agentic Operating-System Operator in Bash,\\ Designed with the PEAS Model}
\newcommand{\coursecode}{CSE000}
\newcommand{\coursetitle}{Course Title}
\newcommand{\proposaldate}{13 November 2026}

% ---------- Team (edit names, IDs, e-mails) ---------------------------------
\newcommand{\memA}{Tamzidul Haque}
\newcommand{\memB}{Umma Iman Monea}
\newcommand{\memC}{Puja Chowdhury}
\newcommand{\memD}{Tamima Rahman}
\newcommand{\memE}{Shagato Chowdhury}

\hypersetup{pdftitle={os-buddy: Project Proposal},pdfsubject={Group project proposal},
  pdfkeywords={agentic AI, Bash, Gemini, PEAS, operating system}}

% ---------- Look and feel (matches the report) ------------------------------
\setlength{\parindent}{1.5em}
\setlength{\parskip}{0pt}
\renewcommand{\arraystretch}{1.15}
\definecolor{hdrgrey}{gray}{0.90}
\newcommand{\code}[1]{\texttt{\small #1}}
\captionsetup{font=small,labelfont=bf,skip=4pt}

\makeatletter
\renewcommand\section{\@startsection{section}{1}{\z@}%
  {1.6ex \@plus .4ex \@minus .2ex}{0.8ex \@plus .2ex}{\normalfont\large\bfseries}}
\renewcommand\subsection{\@startsection{subsection}{2}{\z@}%
  {1.2ex \@plus .3ex \@minus .2ex}{0.5ex \@plus .2ex}{\normalfont\normalsize\bfseries}}
\makeatother

\pagestyle{fancy}
\fancyhf{}
\fancyhead[L]{\footnotesize\projectname: Project Proposal}
\fancyhead[R]{\footnotesize\thepage}
\renewcommand{\headrulewidth}{0.4pt}

% Grey table header and compact lists
\newcommand{\hdrrow}{\rowcolor{hdrgrey}}
\usepackage{colortbl}
\newcommand{\tightlist}{\setlength{\itemsep}{0pt}\setlength{\parsep}{0pt}\setlength{\topsep}{2pt}\setlength{\parskip}{0pt}}
\newcommand{\lead}[1]{\textbf{#1}}
\newcolumntype{R}[1]{>{\raggedright\arraybackslash}p{#1}}

\begin{document}

% =============================================================================
%  COVER PAGE
% =============================================================================
\thispagestyle{empty}
\begin{center}
  \vspace*{1.6cm}
  \includegraphics[width=0.52\textwidth]{diu-logo.png}\\[1.0em]
  {\bfseries Faculty of Science and Information Technology}\\
  Department of Computer Science and Engineering\\[2.4em]
  \rule{\linewidth}{0.5pt}\\[2.2em]
  {\fontsize{40}{44}\selectfont\bfseries \projectname}\\[0.6em]
  {\Large \projectsubtitle}\\[2.2em]
  \rule{\linewidth}{0.5pt}\\[2.2em]
  {\small A project proposal submitted for \quad \textbf{\coursecode\ $|$ \coursetitle}}\\[1.4em]

  {\small\bfseries SUBMITTED BY}\\[0.3em]
  \begingroup\small
  \begin{tabular}{@{}p{4.4cm}p{2.2cm}p{6.0cm}@{}}
    \hline
    \hdrrow \textbf{Name} & \textbf{Student ID} & \textbf{Email} \\ \hline
    \memA & 241-15-877 & {\scriptsize\texttt{tamzidul241-15-877@diu.edu.bd}} \\ \hline
    \memB & 241-15-540 & {\scriptsize\texttt{iman241-15-540@diu.edu.bd}} \\ \hline
    \memC & 241-15-493 & {\scriptsize\texttt{puja241-15-493@diu.edu.bd}} \\ \hline
    \memD & 241-15-030 & {\scriptsize\texttt{rahman241-15-030@diu.edu.bd}} \\ \hline
    \memE & 241-15-331 & {\scriptsize\texttt{shagato241-15-331@diu.edu.bd}} \\ \hline
  \end{tabular}\\[1.2em]

  {\small\bfseries SUPERVISED BY}\\[0.3em]
  \begin{tabular}{@{}p{13.45cm}@{}}
    \hline
    \hdrrow \textbf{Dr. Abdus Sattar} \\ \hline
    Associate Professor and Director, M.Sc. in CSE\\
    Department of Computer Science and Engineering\\
    Faculty of Science and Information Technology, Daffodil International University \\ \hline
  \end{tabular}\\[1.2em]
  \endgroup
  {\small \proposaldate}
\end{center}
\newpage
\setcounter{page}{1}

% =============================================================================
%  BODY (pages 1--2)
% =============================================================================
{\LARGE\bfseries Project Proposal}\\[-0.1em]
\rule{\linewidth}{0.8pt}\vspace{0.4em}

\section{Project Title and Description}

\subsection{Title}
\textbf{\projectname: An Agentic Operating-System Operator in Bash, Designed with the PEAS Model}

\subsection{Description}
Operating a Unix system means remembering many commands, their flags and their risks.
\lead{\projectname} lowers that barrier. The user states a goal in plain English, such as
``find the three largest files in this folder''; the agent plans, runs commands, reads and
edits files, checks the result of every action and continues until the goal is met or it can
report honestly that it has failed. It is written entirely in Bash and uses Google's Gemini
function-calling interface to choose each next step, with the conversation history serving as
its internal state. This makes it a goal-based, model-based agent.

\lead{A fixed safety policy sits between the model and the system.} The policy is Bash code
that the model cannot override: read-only commands run automatically, changes ask for the
user's approval, and dangerous patterns such as \code{rm -rf /} are always blocked. The tool
runs on Linux and Unix-like systems (Bash 4.4+, \code{curl}, \code{jq}) in two modes, an
interactive chat with memory and a one-shot mode (\code{os-buddy "goal"}) with meaningful exit
codes. Two extra commands, \code{os-buddy disk} and \code{os-buddy health}, print system
reports without an API key.

\subsection{Problem and Objectives}
Language models can suggest shell commands, but running them blindly is unsafe and chat-only
assistants leave verification to the user. The project builds an agent that finishes
multi-step tasks while keeping the user in control. Its objectives are:
\begin{itemize}\tightlist
  \item a working plan--act--observe--verify agent loop in pure Bash;
  \item a three-level safety policy (allow, ask, block) with zero bypasses;
  \item a robust Gemini integration (retries, blocked and empty responses) that never
        exposes the API key in process listings or child environments;
  \item an offline test suite and continuous integration on macOS and several Linux versions.
\end{itemize}

\section{Design Approach and Scope}
The agent is specified with the PEAS framework (Table~\ref{tab:peas}). The code is a small
entry point plus focused modules (\code{lib/agent.sh}, \code{gemini.sh}, \code{policy.sh},
\code{tools.sh}, \code{ui.sh}), with prompt text and tool declarations kept as data files.
\emph{Out of scope:} Windows support, a graphical interface, and state-changing actions
without approval.

\begin{table}[h]
\centering
\caption{PEAS summary of \projectname.}
\label{tab:peas}
\small
\begin{tabularx}{\linewidth}{@{}|>{\bfseries}p{2.2cm}|X|@{}}
\hline
\hdrrow Component & \textbf{Specification} \\ \hline
Performance & Goal completion with verification; no policy bypass; recovery from failed
  commands; at most 30 steps; resilience to API errors; exit codes 0 done, 1 error, 2 usage,
  3 step limit. \\ \hline
Environment & Local Linux or Unix host, a human at a terminal (or none), the Gemini API and
  untrusted file content; partially observable, stochastic and sequential. \\ \hline
Actuators & Shell commands, file read and write, user questions and approval prompts, all
  passed through the policy layer. \\ \hline
Sensors & Command output and exit status, file contents, system facts and user replies. \\ \hline
\end{tabularx}
\end{table}

\newpage
\section{Team Contributions}
The project is carried out by a group of five. Work is divided by component rather than by
feature, so each member owns an interface and can work without waiting for the others.

\begin{table}[h]
\centering
\caption{Division of work within the group.}
\label{tab:team}
\small
\begin{tabularx}{\linewidth}{@{}|R{3.3cm}|R{3.6cm}|X|@{}}
\hline
\hdrrow \textbf{Member} & \textbf{Primary responsibility} & \textbf{Principal contribution} \\ \hline
\memE & Architecture, agent core & PEAS design, agent loop, stuck detection, history trimming,
  CLI entry point, chat and one-shot modes, release \\ \hline
\memA & Safety policy & Allow/ask/block classification, dangerous-pattern list, approval
  prompts, secret-scanning commit hook \\ \hline
\memB & Gemini integration & API client and retries, tool declarations, system prompt,
  configuration precedence, API-key protection \\ \hline
\memC & Tools and reports & Tool implementations with timeouts and truncation, \code{disk}
  and \code{health} reports, terminal rendering \\ \hline
\memD & Testing, CI, documentation & Offline test runner and scripted fixtures, CI and ShellCheck,
  README, design document, final report \\ \hline
\multicolumn{3}{|l|}{\cellcolor{hdrgrey}\textbf{Shared by the whole group}} \\ \hline
\multicolumn{3}{|l|}{Requirements analysis, architecture decisions, code review, demonstration} \\ \hline
\end{tabularx}
\end{table}

\section{Plan and Expected Outcome}
\begin{table}[h]
\centering
\caption{Planned schedule.}
\label{tab:plan}
\small
\begin{tabularx}{\linewidth}{@{}|R{1.6cm}|X|R{5.0cm}|@{}}
\hline
\hdrrow \textbf{Weeks} & \textbf{Milestone} & \textbf{Lead} \\ \hline
1--2 & Requirements, PEAS design, repository and tooling & \memE, \memD \\ \hline
3--4 & Gemini client, tool declarations, basic agent loop & \memB, \memE \\ \hline
5--6 & Safety policy, approvals, tools and reports & \memA, \memC \\ \hline
7--8 & Integration, tests and CI, final report, demonstration & \memD, all \\ \hline
\end{tabularx}
\end{table}

\lead{Expected outcome.} A working \code{os-buddy} command with documentation, an offline test
suite and a demonstration of the agent completing real system tasks safely. Each performance
measure has a check: scripted multi-step scenarios for goal completion, a table of bypass
attempts for the policy, stubbed API errors for resilience, and exit-code tests. The main risks
are a harmful command proposed by the model (mitigated by the policy layer), prompt injection
through file content (tool output is treated as data and every change needs approval), and
API outages (retries and clear errors).

\section*{References}
\begingroup\small
\begin{enumerate}\tightlist
  \item S.~Russell and P.~Norvig, \emph{Artificial Intelligence: A Modern Approach}, 4th ed., Pearson, 2020.
  \item Google AI for Developers, ``Gemini API: Function calling,'' \texttt{ai.google.dev}.
  \item GNU Project, \emph{Bash Reference Manual}, \texttt{gnu.org/software/bash}.
\end{enumerate}
\endgroup

\end{document}
